\documentclass[a4paper,11pt]{article}
\usepackage[utf8]{inputenc}
\usepackage[T1]{fontenc}
\usepackage[french]{babel}
\usepackage{amsmath,amsfonts,amssymb}
\usepackage{geometry}
\usepackage{fancyhdr}
\usepackage{titlesec}
\usepackage{enumitem}

\geometry{margin=2.5cm}

% Header and footer configuration
\pagestyle{fancy}
\fancyhf{}
\fancyhead[L]{Licence 1 Informatique}
\fancyhead[R]{CERI - Avignon Université}
\fancyfoot[C]{Mathématiques Discrètes \hfill TD6 - Application de la combinatoire \hfill Page \thepage/4}

% Title formatting
\titleformat{\section}
{\Large\bfseries}
{\thesection}
{1em}
{}

\titleformat{\subsection}
{\large\bfseries}
{\thesubsection}
{1em}
{}

\begin{document}

\begin{center}
{\Huge\bfseries TD6 - Application de la combinatoire} \\[0.5cm]
{\large Hiver 2026} \\[0.3cm]
\end{center}

\vspace{0.5cm}

L'objectif de ce TD est de mettre en pratique les concepts de la combinatoire tels que le principe des tiroirs, les permutations, combinaisons et d'appliquer les connaissances sur le coefficient binomial.

\section{Quelques calculs}

\begin{enumerate}
\item Un tiroir contient 24 chaussettes en vrac, la moitié sont bleues, la moitié sont noires. On en choisit au hasard sans regarder.
\begin{enumerate}
\item Combien doit-on en tirer au sort pour être sûr d'en avoir une paire de la même couleur ?
\item Combien doit-on en tirer au sort pour être sûr d'en avoir une paire noire ?
\end{enumerate}

\item Un saladier contient 10 boules rouges et 10 boules bleues. On en choisit au hasard sans regarder.
\begin{enumerate}
\item Combien doit-on en tirer au sort pour être sûr d'en avoir trois de la même couleur ?
\item Combien doit-on en tirer au sort pour être sûr d'en avoir trois rouges ?
\end{enumerate}

\item Un lycée compte 1500 élèves.
\begin{enumerate}
\item Montrez qu'il existe au moins trois élèves qui ont les mêmes initiales.
\item Quel est le plus petit nombre d'élèves nécessaires pour garantir qu'au moins trois ont les mêmes initiales ?
\end{enumerate}

\item On considère les étudiants d'une université qui sont nés dans un département français. On rappelle qu'il y a 101 départements français.
\begin{enumerate}
\item Combien d'étudiants nés en France doivent être inscrits à l'université pour être sûr qu'il y en a au moins 40 nés dans le même département ?
\item Combien d'étudiants nés en France doivent être inscrits à l'université pour être sûr qu'il y en a au moins 10 nés dans le Vaucluse ?
\end{enumerate}

\item Une assemblée réservée aux étudiants se tient à l'université dans un amphi qui comporte 800 places et qui est plein.
\begin{enumerate}
\item Montrez qu'il existe un niveau de diplôme (licence, master ou doctorat) représenté par au moins $x$ étudiants et précisez cette valeur $x$.
\item Après sondage, on sait que l'amphi comporte 133 étudiants inscrits en doctorat. Montrez qu'au moins 134 étudiants sont inscrits dans la même année de licence ou master.
\end{enumerate}
\end{enumerate}

\section{Diviseurs}

\begin{enumerate}
\item \emph{On considère un ensemble} $E \subset \mathbb{N}$ \emph{avec} $|E| = 3$. Montrez qu'il existe deux éléments pairs ou deux éléments impairs dans $E$.

\item \emph{On considère un ensemble} $E \subset \mathbb{N}$ \emph{avec} $|E| = 5$. Montrez qu'il existe deux éléments de $E$ qui ont le même reste dans la division entière par 4.

\item \emph{On considère un ensemble} $E \subset \mathbb{N}$ \emph{avec} $|E| = d + 1$, avec $d \in \mathbb{N}^*$. Montrez qu'il existe deux éléments de $E$ qui ont le même reste dans la division entière par $d$.

\item Soit $n \in \mathbb{N}^*$. \emph{On considère une liste de} $n$ \emph{entiers consécutifs, montrez qu'il y en a exactement un qui est divisible par} $n$.
\end{enumerate}

\section{Sommes}

\begin{enumerate}
\item On considère l'ensemble $E = \{1, 2, 3, 4, 5, 6, 7, 8\}$.
\begin{enumerate}
\item Est-il vrai que lorsqu'on tire au hasard 5 éléments dans $E$, il y en a deux parmi eux dont la somme fait 9 ?
\item Est-il vrai que lorsqu'on tire au hasard 4 éléments dans $E$, il y en a deux parmi eux dont la somme fait 9 ?
\end{enumerate}

\item On considère l'ensemble $E = \{1, 2, 3, 4, 5, 6, 7, 8, 9, 10\}$.
\begin{enumerate}
\item Est-il vrai que lorsqu'on tire au hasard 7 éléments dans $E$, il y aura deux paires parmi eux dont la somme fait 11 ?
\item Est-il vrai que lorsqu'on tire au hasard 6 éléments dans $E$, il y aura deux paires parmi eux dont la somme fait 11 ?
\end{enumerate}

\item On considère l'ensemble $E = \{1, 2, 3, 4, 5, 6\}$.
\begin{enumerate}
\item Combien d'éléments doit-on tirer dans $E$ pour être certain d'avoir au moins une paire dont la somme fait 7 ?
\item Combien d'éléments doit-on tirer dans $E$ pour être certain d'avoir au moins deux paires dont la somme fait 7 ?
\end{enumerate}

\item On considère l'ensemble $E = \{1, 3, 5, 7, 9, 11, 13, 15\}$.
\begin{enumerate}
\item Combien d'éléments doit-on tirer dans $E$ pour être certain d'avoir au moins une paire dont la somme fait 16 ?
\item Combien d'éléments doit-on tirer dans $E$ pour être certain d'avoir au moins deux paires dont la somme fait 16 ?
\end{enumerate}
\end{enumerate}

\section{Tournois}

\begin{enumerate}
\item Pendant un mois de 30 jours, une équipe de football joue au moins un match par jour, et joue au plus 45 matches au total. Montrez qu'il y a une période d'un certain nombre de jours consécutifs au cours desquels l'équipe aura disputé 14 matches.

\item Un boxeur est resté champion pendant une période de 75 semaines consécutives. Il a fait au moins un combat par semaine, et a fait au plus 125 combats en tout. Montrez qu'il y a une période d'un certain nombre de semaines consécutives au cours desquelles il a fait 24 combats.
\end{enumerate}

\section{Amis ou ennemis}

\begin{enumerate}
\item \emph{On considère un groupe de six personnes dans lequel n'importe quelle paire de deux personnes sont soit amis soit ennemis. Montrez qu'il existe dans le groupe trois personnes amis ou trois personnes ennemies.}
\end{enumerate}

\section{Application des permutations et combinaisons}

\begin{enumerate}
\item Soit l'ensemble $E = \{\spadesuit, \clubsuit, \heartsuit, \diamondsuit\}$
\begin{enumerate}
\item Combien de permutations y a-t-il de l'ensemble $E$ ?
\item Combien $k$-combinaisons y a-t-il de l'ensemble $E$ ? Pour $k = 0, 1, 2, 3, 4$
\item Combien $k$-permutations y a-t-il de l'ensemble $E$ ? Pour $k = 0, 1, 2, 3, 4$
\item Listez toutes les 2-combinaisons de $E$.
\item Listez quelques 2-permutations de $E$.
\end{enumerate}

\item Soit l'ensemble $E = \{1, 2, 3, 4, 5\}$
\begin{enumerate}
\item Combien de permutations y a-t-il de l'ensemble $E$ ?
\item Combien $k$-combinaisons y a-t-il de l'ensemble $E$ ? Pour $k = 0, 1, 2, 3, 4, 5$
\item Combien $k$-permutations y a-t-il de l'ensemble $E$ ? Pour $k = 0, 1, 2, 3, 4, 5$
\item Listez toutes les 3-combinaisons de $E$.
\item Listez quelques 3-permutations de $E$.
\end{enumerate}

\item Soit l'ensemble $E = \{A, B, C, D, E, F\}$
\begin{enumerate}
\item Combien de permutations y a-t-il de l'ensemble $E$ ?
\item Combien $k$-combinaisons y a-t-il de l'ensemble $E$ ? Pour $k = 0, 1, 2, 3, 4, 5, 6$
\item Combien $k$-permutations y a-t-il de l'ensemble $E$ ? Pour $k = 0, 1, 2, 3, 4, 5, 6$
\item Listez toutes les 4-combinaisons de $E$.
\item Listez quelques 4-permutations de $E$.
\end{enumerate}
\end{enumerate}

\section{Vie d'association}

Une association compte 25 membres.

\begin{enumerate}
\item Le bureau est composé de 4 membres. Combien y a-t-il de façons de sélectionner les membres du bureau ?

\item Le bureau est composé d'un président, un vice-président, un secrétaire et un trésorier. Combien y a-t-il de façons de sélectionner les membres du bureau sachant qu'une personne ne peut occuper qu'un seul poste au maximum ?
\end{enumerate}

\section{Sélection d'équipe}

Dans un groupe de 13 personnes, on doit en sélectionner 10 pour former une équipe.

\begin{enumerate}
\item Combien y a-t-il de façons de former cette équipe ?

\item Combien y a-t-il de façons de former cette équipe si chaque rôle dans l'équipe est différent ?

\item Dans le groupe, il y a 3 femmes et 10 hommes. Combien y a-t-il de façons de former l'équipe (avec des rôles identiques) si on souhaite qu'au moins une femme soit sélectionnée ?
\end{enumerate}

\section{Loterie}

Cent tickets de loterie sont distribués à 100 personnes différentes. Il y a 4 prix différents à gagner, dont un voyage.

\begin{enumerate}
\item Combien de façons y a-t-il de désigner les gagnants ?

\item Combien de façons y a-t-il de désigner les gagnants si la personne avec le billet 47 gagne le voyage ?

\item Combien de façons y a-t-il de désigner les gagnants si la personne avec le billet 47 gagne un prix ?

\item Combien de façons y a-t-il de désigner les gagnants si la personne avec le billet 47 ne gagne pas de prix ?

\item Combien de façons y a-t-il de désigner les gagnants si la personne avec le billet 19 et celle avec le billet 47 gagnent un prix ?

\item Combien de façons y a-t-il de désigner les gagnants si la personne avec le billet 73, celle avec le billet 19 et celle avec le billet 47 gagnent un prix ?

\item Combien de façons y a-t-il de désigner les gagnants si la personne avec le billet 97, celle avec le billet 73, celle avec le billet 19 et celle avec le billet 47 gagnent un prix ?

\item Combien de façons y a-t-il de désigner les gagnants si ni la personne avec le billet 97, ni celle avec le billet 73, ni celle avec le billet 19 ni celle avec le billet 47 ne gagnent un prix ?

\item Combien de façons y a-t-il de désigner les gagnants si le voyage a été gagné par la personne avec le billet 97, ou celle avec le billet 73, ou celle avec le billet 19 ou celle avec le billet 47 ?

\item Combien de façons y a-t-il de désigner les gagnants si la personne avec le billet 19 et celle avec le billet 47 gagnent un prix et que la personne avec le billet 73 et celle avec le billet 97 ne gagnent pas de prix ?
\end{enumerate}

\section{File d'attente}

\begin{enumerate}
\item Combien de façon y a-t-il d'ordonner 8 hommes et 5 femmes dans une file d'attente s'il n'y a jamais deux femmes l'une derrière l'autre ?

\item Combien de façon y a-t-il d'ordonner 6 hommes et 10 femmes dans une file d'attente s'il n'y a jamais deux hommes l'un derrière l'autre ?
\end{enumerate}

\section{Exploiter des formules}

\begin{enumerate}
\item Développer la formule $(x + y)^6$

\item Combien y a-t-il de termes dans le développement de $(x + y)^{100}$

\item Donnez le coefficient de $x^9y^4$ dans $(x + y)^{13}$

\item Donnez le coefficient de $x^{12}y^{13}$ dans $(2x - 3y)^{25}$

\item Donnez le coefficient de $x^7$ dans $(1 + x)^{14}$

\item Donnez le coefficient de $x^9$ dans $(2 - x)^{19}$

\item Donnez le coefficient de $x^8y^9$ dans $(3x + 2y)^{17}$

\item Donnez le coefficient de $x^{104}y^{20}$ dans $(2x - 3y)^{20}$
\end{enumerate}

\section{Égalités remarquables}

\begin{enumerate}
\item Démontrez la relation de Pascal que nous rappelons ici : $\binom{n+1}{k} = \binom{n}{k-1} + \binom{n}{k}$

\item Démontrez l'équivalence remarquable suivante : $\binom{n}{r} \cdot \binom{r}{k} = \binom{n}{k} \cdot \binom{n-k}{r-k}$

\item Démontrez l'équivalence remarquable suivante : $\binom{n-1}{k} \cdot \binom{n}{k+1} = \binom{n+1}{k} \cdot \binom{n-1}{k+1}$

\item \emph{Pour aller plus loin : donnez des preuves combinatoires des égalités précédentes}
\end{enumerate}

\end{document}